\begin{Solbook}{1}
 \par \noindent  \begin {enumerate} \item Ta có \begin {align*} &(2x-1)^2=9(x-2)^2\\ &(2x-1)^2-\big [3(x-2)\big ]^2=0\\ &\big [(2x-1)-3(x-2)\big ]\big [(2x-1)+3(x-2)\big ]=0\\ &(5-x)(5x-7)=0. \end {align*} Do đó $x=5$ hoặc $x=\dfrac {7}{5}$. \par \item Thu gọn vế phải rồi nhân cả hai vế với $12>0$: \begin {align*} 3-\dfrac {2x-5}{4}&\le \dfrac {x+5}{6}\\ 3(17-2x)&\le 2(x+5)\\ 51-6x&\le 2x+10\\ 41&\le 8x\\ x&\ge \dfrac {41}{8}. \end {align*} Vậy bất phương trình có nghiệm $x\ge \dfrac {41}{8}$. \par \item Dùng phương pháp cộng đại số, ta có \begin {align*} &\begin {cases} \dfrac {3x-y}{4}+\dfrac {x+2y}{3}=\dfrac {1}{6},\\ 2(x+1)-3(y-2)=-3 \end {cases}\\[0.7em] \Leftrightarrow \quad &\begin {cases} 13x+5y=2,\\ 2x-3y=-11 \end {cases}\\[0.7em] \Leftrightarrow \quad &\begin {cases} 39x+15y=6,\\ 10x-15y=-55 \end {cases}\\[0.7em] \Leftrightarrow \quad &\begin {cases} 49x=-49,\\ 2x-3y=-11 \end {cases}\\[0.7em] \Leftrightarrow \quad &\begin {cases} x=-1,\\ y=3. \end {cases} \end {align*} Vậy hệ phương trình có nghiệm $(x;y)=(-1;3)$. \end {enumerate}  \hfill \qedEX 
\end{Solbook}
\begin{Solbook}{2}
 \par \noindent  \noindent a) $A=4$ \par \par \noindent b) $B=\dfrac {5-\sqrt {x}}{\sqrt {x}+5}$  \hfill \qedEX 
\end{Solbook}
\begin{Solbook}{3}
 \par \noindent  Một cái bánh mỳ có giá 20\,000 đồng, một cái bánh bao có giá 25\,000 đồng.  \hfill \qedEX 
\end{Solbook}
\begin{Solbook}{4}
 \par \noindent Khoảng $35{,}7$ km. \hfill \qedEX 
\end{Solbook}
\begin{Solbook}{5}
 \par \noindent \immini {\begin {enumerate} \item $MN$ là đường trung trực của $SO$ \item $\triangle SNF \backsim \triangle SEN$ suy ra $SN^2=SF\cdot SE$ \par $\triangle SNH \backsim \triangle SON$ suy ra $SN^2=SH\cdot SO$ \item $\triangle SHI \backsim \triangle SKO$ suy ra $SK\cdot SI=SH\cdot SO$ \par Chứng minh $SF+SE=2SK$ \par Suy ra $SI(SF+SE)=2SI\cdot SK=2SE\cdot SF$ \end {enumerate}} {\begin {tikzpicture}[scale=1, font=\footnotesize ] \coordinate (O) at (0,0); \def \R {2} \draw (O) circle(\R ); \coordinate (S) at (-5,0); \path [name path=circleO] (O) circle(\R ); \path [name path=circleS] (-2.5,0) circle(2.5); \path [name intersections={of=circleS and circleO, by={M,N}}]; \coordinate (E) at ($-1*(N)$); \path [name path=os] (O)--(S); \path [name path=mn] (M)--(N); \path [name intersections={of=os and mn, by=H}]; \path [name path=se] (S)--(E); \path [name intersections={of=circleO and se, by={F1,F2}}]; \coordinate (F) at (F2); \path [name intersections={of=se and mn, by=I}]; \draw (S)--(M)--(O)--cycle; \draw (M)--(N)--(E)--(S); \draw (S)--(O); \coordinate (K) at ($(F)!0.5!(E)$); \draw (O)--(K); \foreach \p /\pos in {S/left,M/above,N/below,O/below right,F/below,E/above,K/above,I/below left,H/below left} {\draw [fill=white] (\p ) circle(1.5pt) node[\pos ] {$\p $}; } \end {tikzpicture}}  \hfill \qedEX 
\end{Solbook}
\begin{Solbook}{6}
 \par \noindent  $x=1$ \par Gợi ý: cộng hai vế của phương trình với -2: \[\dfrac {x-1}{2022}+\dfrac {x-1}{2023}>\dfrac {x-1}{2024}+\dfrac {x-1}{2025}\]  \hfill \qedEX 
\end{Solbook}
